\documentclass[11pt]{article}
\usepackage[margin=1in]{geometry}
\usepackage[T1]{fontenc}
\usepackage[utf8]{inputenc}
\usepackage{amsmath,amssymb,amsthm}
\usepackage{enumitem}

\title{ICPC World Finals 2020\\C. Domes}
\author{}
\date{}

\begin{document}
\maketitle

\section*{Problem Summary}

We are given $n$ dome positions inside the rectangle $[0,dx] \times [0,dy]$ and a permutation
$p_1, p_2, \dots, p_n$.
We must find the area of all photographer positions $X$ inside the rectangle from which one can take a
180-degree photo whose left-to-right order of domes is exactly
\[
    p_1, p_2, \dots, p_n.
\]

\section*{Key Observations}

\begin{itemize}[leftmargin=*]
    \item In a 180-degree photo, the domes are ordered exactly by the angular order of the rays from the photographer.
    Therefore for every pair $i < j$, dome $p_i$ must appear before dome $p_j$ in that angular order.
    \item For a fixed pair $i < j$, the condition ``$p_i$ is before $p_j$'' is equivalent to saying that the
    photographer lies on the \emph{right-hand side} of the directed line from dome $p_i$ to dome $p_j$.
    In coordinates, if $A$ is the position of $p_i$, $B$ the position of $p_j$, and $X$ the photographer, then we need
    \[
        \operatorname{cross}(B-A,\; X-A) \le 0.
    \]
    \item Hence the feasible region is the intersection of:
    \begin{itemize}[leftmargin=*]
        \item the original rectangle,
        \item one half-plane for every pair $i < j$.
    \end{itemize}
    \item An intersection of half-planes is convex, so we can start from the rectangle and clip it repeatedly.
    Since $n \le 100$, there are at most $\binom{100}{2} = 4950$ half-planes, which is small enough for ordinary
    polygon clipping.
\end{itemize}

\section*{Algorithm}

\begin{enumerate}[leftmargin=*]
    \item Reorder the dome coordinates according to the requested permutation
    $p_1, p_2, \dots, p_n$.
    \item Initialize the current polygon as the rectangle
    \[
        (0,0),\ (dx,0),\ (dx,dy),\ (0,dy).
    \]
    \item For every pair $i < j$:
    \begin{itemize}[leftmargin=*]
        \item let $A$ be dome $p_i$ and $B$ be dome $p_j$;
        \item clip the current polygon with the half-plane
        \[
            \operatorname{cross}(B-A,\; X-A) \le 0.
        \]
    \end{itemize}
    We use the standard Sutherland--Hodgman clipping step for one convex polygon against one line.
    \item After all clips, compute the area of the remaining polygon with the shoelace formula.
    If the polygon becomes empty, the answer is $0$.
\end{enumerate}

\section*{Correctness Proof}

We prove that the algorithm returns the correct answer.

\paragraph{Lemma 1.}
For two domes $A$ and $B$, with $A$ required to appear before $B$ in the photo, a photographer position $X$ is valid
for this pair if and only if $X$ lies on the right-hand side of the directed line $A \to B$.

\paragraph{Proof.}
The left-to-right order in a 180-degree photo is the angular order of the rays from $X$ to the domes inside one
visible semicircle.
So ``$A$ appears before $B$'' means that when rotating from ray $XA$ to ray $XB$, we move clockwise by an angle strictly
less than $\pi$.
That is exactly the condition that the oriented triangle $(A,B,X)$ has non-positive signed area, or equivalently
\[
    \operatorname{cross}(B-A,\; X-A) \le 0.
\]
The boundary case is collinearity; it has area $0$, so including it does not affect the final answer. \qed

\paragraph{Lemma 2.}
If a position $X$ satisfies the half-plane condition for every pair $i < j$, then there exists a 180-degree photo from
which the domes appear in the requested order $p_1, p_2, \dots, p_n$.

\paragraph{Proof.}
Let $\alpha_k$ be the polar angle of the ray from $X$ to dome $p_k$.
For every $i < j$, Lemma 1 gives that the directed angle from ray $Xp_i$ to ray $Xp_j$ is clockwise and smaller than
$\pi$.
Therefore the angles can be lifted to real numbers so that
\[
    \alpha_1 > \alpha_2 > \dots > \alpha_n
    \quad\text{and}\quad
    \alpha_1 - \alpha_n < \pi.
\]
So all rays lie inside one open semicircle, and their angular order inside that semicircle is exactly
$p_1, p_2, \dots, p_n$.
Point the camera toward the middle direction of that semicircle.
Then all domes are visible and appear in the required left-to-right order. \qed

\paragraph{Lemma 3.}
If there exists a valid photo from position $X$ whose dome order is $p_1, p_2, \dots, p_n$, then $X$ satisfies all the
half-plane constraints used by the algorithm.

\paragraph{Proof.}
In a valid photo, for every pair $i < j$, dome $p_i$ appears before dome $p_j$.
By Lemma 1, this implies that $X$ lies on the right-hand side of the directed line from $p_i$ to $p_j$.
Hence all pair constraints hold. \qed

\paragraph{Theorem.}
The algorithm outputs the correct answer for every valid input.

\paragraph{Proof.}
By Lemma 3, every valid photographer position lies in every half-plane used by the algorithm, so it survives all
clipping steps.
By Lemma 2, every position that survives all clipping steps is indeed valid for the requested photo order.
Thus the final polygon produced by the algorithm is exactly the feasible region inside the rectangle.
The shoelace formula therefore computes exactly the desired area. \qed

\section*{Complexity Analysis}

There are $O(n^2)$ half-planes.
Each clipping step is linear in the current polygon size, which is also $O(n^2)$ in the worst case.
Therefore the running time is $O(n^4)$ in the worst case, which is easily fast enough for $n \le 100$.

The polygon itself has $O(n^2)$ vertices, so the memory usage is $O(n^2)$.

\section*{Implementation Notes}

\begin{itemize}[leftmargin=*]
    \item The algorithm clips with \emph{all} pairs $i < j$, not only consecutive pairs in the permutation.
    \item Using \texttt{long double} for the geometric computation is sufficient for the required $10^{-3}$ accuracy.
    \item After clipping, removing consecutive duplicate vertices helps keep the polygon stable numerically.
\end{itemize}

\end{document}
